\section{Symmetries and impartial games}
\label{sec:actions-games}
The seed article connects combinatorics with group theory and lists
combinatorial game theory among its related fields.
The two applications below use different finite computations. Symmetry
reduces a configuration to its least representative; a game assigns its
positions numbers in order of increasing rank.

\subsection{Finite group actions and color inventories}
A finite group is specified by a size $s>0$, a multiplication table
on $[s]$, an identity $e$ and an inverse table. Check the identity,
inverse and associative laws by loops of length at most $s^3$.
A left action on $X=[u]$ is a table with
$ex=x$ and $(gh)x=g(hx)$.
Write $Gx=\{gx:g\in G\}$, $G_x=\{g:gx=x\}$,
$X^g=\{x:gx=x\}$, and $X/G$ for the set of least members of the
orbits. Empty $X$ gives no representatives.
\begin{theorem}[Orbit--stabilizer and the orbit-counting lemma]\label{result63}
For every $x\in X$,
\[
 |Gx|\,|G_x|=s,\qquad
 s\,|X/G|=\sum_{g\in G}|X^g|.
\]
In particular the size of each orbit divides $s$ and the sum of the
numbers of fixed points is divisible by $s$.
\end{theorem}
\begin{proof}
Inverse elements make the relation $y=gx$ symmetric, while the identity
and multiplication give reflexivity and transitivity. For each $y\in Gx$
take the least $t_y$ satisfying $t_yx=y$. The map
$h\mapsto t_yh$ from $G_x$ onto the elements sending $x$ to $y$ has
inverse $g\mapsto t_y^{-1}g$. Summing these fiber sizes gives the first
formula. Moreover $h\mapsto t_yht_y^{-1}$ carries $G_x$ bijectively onto
$G_y$; conjugate by $t_y^{-1}$ for its inverse. Count pairs $(g,x)$
with $gx=x$. Each orbit contributes $|Gx|\,|G_x|=s$, while a fixed
$g$ contributes $|X^g|$. This proves the second identity.
\end{proof}
If $H$ is a subgroup of $G$, apply the first identity to left
multiplication by $H$ on each left-translation orbit $Hx$. It has
$|H|$ elements since $hx=x$ implies $h=e$. The disjoint orbits partition
$G$, giving the finite Lagrange divisibility statement $|H|\mid s$.

\begin{theorem}[P\'olya's inventory formula]\label{result64}
Let $G$ act on $m$ places. With $q$ available colors, give each coloring
$w$ the monomial $\prod_{i<m} z_{w(i)}$. Let $Z$ be the sum of one such
monomial per orbit of colorings. If $c_j(g)$ is the number of $j$-cycles
of $g$ on the places, then
\[
 |G| Z=\sum_{g\in G}\prod_{j=1}^{m}
          (z_0^j+\cdots+z_{q-1}^j)^{c_j(g)}.
\]
Taking all weights equal to one yields
$|G|\,\#\{\text{coloring orbits}\}=\sum_g q^{\sum_j c_j(g)}$.
\end{theorem}
\begin{proof}
The induced action is $(gw)(i)=w(g^{-1}i)$. Its fixed colorings are
exactly those constant on each cycle of $g$. Starting at the least
unvisited place, iterate $g$ until it returns. Injectivity forces the
first repeated place to be the start, and the return takes at most $m$
steps. Repeat while an unvisited place remains. A cycle of length $j$
independently receives any color $a$ and contributes $z_a^j$.
Color counts are constant on an orbit. Apply the orbit-counting
identity separately to the finite invariant set for each inventory
$(a_0,\ldots,a_{q-1})$, $\sum a_i=m$. This proves the coefficient
identity and covers $m=0$ by the empty product convention.
\end{proof}

\subsection{Necklaces, primitive words, and bracelets}
A necklace is a word of length $n\ge1$ modulo cyclic rotations of positions.
Let $N_q(n)$ count necklaces on $q$ symbols; let $P_q(n)$ count primitive
necklaces, whose words have no nonidentity rotation fixing them. Write
$\varphi(d)$ for the number of $a\in[d]$ coprime to $d$, including
$\varphi(1)=1$.
\begin{theorem}[Necklace and primitive necklace formulas]\label{result65}
With $\mu(d)=0$ if a prime square divides $d$ and
$\mu(d)=(-1)^r$ if $d$ is the product of $r$ distinct primes, one has
\[
 n N_q(n)=\sum_{d\mid n}\varphi(d)q^{n/d},
 \qquad
 n P_q(n)=\sum_{d\mid n}\mu(d)q^{n/d}.
\]
Every primitive necklace has a unique lexicographically smallest rotation,
a word strictly smaller than all its nontrivial rotations.
\end{theorem}
\begin{proof}
A rotation by $a\in[n]$ generates the subsets of residues differing by a
multiple of $g=\gcd(n,a)$. Indeed its least positive return time at any
position is $n/g$: writing $a=ga'$, $n=gn'$, the divisibility $n\mid ta$
reduces to $n'\mid ta'$ and hence to $n'\mid t$. The cancellation follows
from the integer B\'ezout identity for the coprime pair, obtained by
recording the linear combination at each Euclidean division; each
nonterminal remainder strictly decreases. This also covers $a=0$.
There are $g$ cycles, so $q^g$ words are fixed. The rotations with
$n/g=d$ correspond to $a=(n/d)b$ with $b\in[d]$, $\gcd(b,d)=1$.
The orbit-counting lemma now gives the first formula.

For an individual word choose the least positive rotational period $t\le n$,
allowing $n$ itself. Periods taken modulo $n$ are closed under addition
and subtraction. Divide $n$ by $t$; its remainder is a period smaller than
$t$ and therefore is zero. Thus $t\mid n$ and the word is a repetition
of its length-$t$ prefix. That prefix is primitive: a smaller rotational
period of it would extend to the whole word. Conversely repeating a
primitive word of length $t\mid n$ gives least period $t$, as any shorter
period would also rotate the length-$t$ prefix into itself (use its
periodic extension). Let $A_t$ count primitive \emph{words}. We have
$q^n=\sum_{t\mid n}A_t$ and $A_n=nP_q(n)$, since each primitive orbit
has $n$ different rotations. For $v>1$, square-free divisors of $v$
are subsets of its distinct prime divisors; toggling the least prime
pairs their signs. Hence $\sum_{d\mid v}\mu(d)=[v=1]$, including $v=1$.
Substitute the preceding expression for $q^{n/d}$ into the signed sum
and regroup the terms with fixed $t$; the multiplier is
$\sum_{d\mid n/t}\mu(d)$, which leaves exactly $A_n$. The prime list here
can be obtained by repeatedly removing the least divisor greater than
one. It is prime by minimality; Euclid's prime-divisor lemma follows from
the same B\'ezout cancellation, giving uniqueness of multiplicities by
successive cancellation. All searches and repeated divisions are bounded
by $n$. Finally take the least in the finite lexicographic list of the
$n$ distinct rotations of a primitive word.
\end{proof}

Let $J_q(n)$ count bracelets, obtained by identifying both rotations and
reversals. These are orbits of the $2n$ formal symmetries
$i\mapsto i+a$ and $i\mapsto a-i\pmod n$. Using formal pairs (sign,
offset) keeps the group of order $2n$ even when two induced permutations
coincide. A reflection has $f$ fixed places determined by $2i=a\pmod n$,
and all remaining places form pairs. For odd $n$ there is precisely one
solution: multiplication by $2$ is invertible modulo $n$. For even $n$,
an odd $a$ has none, and an even $a$ has the two solutions $a/2$ and
$a/2+n/2$. The orbit-counting formula therefore gives
\[
 2n J_q(n)= n N_q(n)+
 \begin{cases}
 n q^{(n+1)/2},& n\text{ odd},\\
 (n/2)(q^{n/2+1}+q^{n/2}),& n\text{ even}.
 \end{cases}
\]
For example, the six binary places have $14$ necklaces, $9$ primitive
necklaces, and $13$ bracelets, obtained by substituting $q=2,n=6$.

\subsection{Winning moves from a finite game table}
Consider finite, impartial, perfect-information games in normal play.
A position has the same legal successors for either player. A player who
cannot move loses. Supply the adjacency table on $[v]$ and a natural
rank $r_i$ for each position, testing $i\to j\Rightarrow r_j<r_i$.
For a finite set $S$ of naturals, $\operatorname{mex} S$ is the least
natural missing from $S$. If $S$ consists of at most $a$ values, some
value in $0,\ldots,a$ is missing by counting, so the search has an explicit
bound. Define the Grundy value $g(i)$ in increasing order of rank by
\[
 g(i)=\operatorname{mex}\{g(j):i\to j\}.
\]
In particular $g(i)\le\#\{j:i\to j\}\le v$. Ties in rank may be broken
by index.

\begin{theorem}[Normal play, the sum theorem, and Nim]\label{result66}
A current player has a winning strategy exactly when $g(i)\ne0$.
In the disjunctive sum, where exactly one component changes at each move,
\[
 g(i,i')=g(i)\mathbin{\mathrm{xor}}g(i').
\]
A Nim heap of size $h$ has value $h$, and a finite list of heaps is
losing exactly when their bitwise exclusive-or is zero. In particular,
two equal heaps are losing. Every finite ranked impartial position is
indistinguishable in win/loss outcome, under addition of any other such
position, from the Nim heap of its value.
\end{theorem}
\begin{proof}
By definition of mex every value strictly below $g(i)$ is realized by
a successor, and $g(i)$ itself is absent. A zero position has only
nonzero successors. From a nonzero position take the least-index successor
of value zero. Starting from a nonzero position, this rule restores zero
after every move by the player following it. Ranks decrease, so the
opponent eventually encounters a terminal position. From a zero position
the opponent can use this strategy after any first move (or win immediately
if no move is available). This establishes both directions and specifies
the strategy table.

Write $a=g(i)$, $b=g(i')$, and $z=a\mathbin{\mathrm{xor}}b$. Induct on
the sum of the two ranks, tabulating all pairs up to the given sum.
A move in the first component yields $a'\mathbin{\mathrm{xor}}b$ with
$a'\ne a$, hence cannot yield $z$ by bitwise cancellation. A move in the
second component also cannot yield $z$. To see every $t<z$, find the
highest differing bit $h$ of $t$ and $z$. At $h$, $z$ has one and $t$
zero, so precisely one of $a,b$ has one there. If it is $a$, put
$a'=t\mathbin{\mathrm{xor}}b$. Above $h$, $a'$ agrees with $a$; at $h$
it is zero where $a$ is one. Thus $a'<a$ and is realized by a successor.
The resulting sum value is $t$. If $b$ has the one, interchange the
components. These facts characterize mex even for $z=0$. Binary digits and
highest differing positions can be scanned up to $1+v+v'$ for the
corresponding table sizes.

The available moves of a heap $h$ go to exactly $0,\ldots,h-1$; finite
induction gives its value $h$. Applying the two-component result repeatedly
and using bitwise associativity gives the assertion for a finite list.
The value determines all joint outcomes with a further component of value
$c$, since the joint value vanishes exactly when $c$ equals that value.
Two different values can be distinguished by adjoining a heap of the
first value.
\end{proof}
For an input giving just a finite acyclic graph, a suitable ranking is
itself primitive recursive. Repeatedly remove the least vertex with no
successor among the remaining vertices, in at most $v$ rounds, assigning
ranks in removal order. If a nonempty remainder had no such vertex,
following least successors $v$ times from a remaining vertex would give
$v+1$ visits and a repeated vertex, hence a directed cycle. Thus all
vertices are removed and every move lowers the resulting rank.
For the partizan version with two move tables and the same rank condition,
tabulate two Boolean entries per position, one for each player to move.
Each entry is true exactly when some legal move leaves a false entry for
the other player; a terminal entry is false. Induction on rank gives the
winner and bounded search selects a winning strategy whenever the entry
is true. The mex and exclusive-or statements above use impartiality.

\paragraph{Lean implementation.}
Represent multiplication, inverses and actions by bounded digit tables
and verify each law by a Boolean fold. Compute \texttt{orbRep x} as
the minimum of the $s$ images, and \texttt{trans x y} as the first $g$
with $gx=y$ under the orbit guard. Extend selectors by zero outside their
guards. Use the two inverse maps in the proof to instantiate finite
bijection counting. The inventory calculation uses
bounded vectors in $[m+1]^q$: count each such vector by scanning at most
$\max(2,q)^m$ colorings. For words, store all positions of \texttt{rotate}
and \texttt{reverseRotate} modulo $n$ and test equality digitwise; verify
cyclic least-period claims by Euclidean division and a single scan of the
positions. In the game computation, store each mex table as $v$ digits
in base $v+1$ (use base two for $v\le1$). Compute new entries in the
specified topological order and prove the successor equation simultaneously
for completed entries by prefix induction. For a product use $vv'$ slots
and rank sums. Turn the exclusive-or argument into two bit lemmas:
involution after flipping the same mask twice, and strict comparison at
the greatest differing bit. Both follow by expanding the binary words
above, at, and below that bit, since the sum of all lower place values
is less than the place value at the bit. Every assertion about all
positions in a completed prefix or rank layer is a bounded Boolean fold.
